\documentclass[a4paper,fontset=none]{ctexart}

\usepackage{indentfirst}
\setlength{\parindent}{0em}
\usepackage{autobreak}
\allowdisplaybreaks
\usepackage{parskip}
\usepackage{geometry}
\geometry{top=3cm,bottom=3cm,left=2cm,right=2cm}
\usepackage{anyfontsize}

\usepackage{fontspec}
\setmainfont{STIX Two Text}
\usepackage{unicode-math}
\unimathsetup{mathrm=sym,mathbf=sym,mathsf=sym,bold-style=ISO}
\setmathfont{STIX Two Math}
\usepackage{physics2}
\usephysicsmodule{ab,op.legacy}
\usepackage{fixdif,derivative}
\newcommand{\um}{\upmu\mathrm{m}}
\newcommand{\ang}{\text{\r{A}}}

\setCJKmainfont{Source Han Serif CN}[BoldFont=Source Han Sans CN Medium,ItalicFont=FZKai-Z03]

\title{\textbf{星际介质天文学作业}}
\author{张植竣}
\date{2024年3月8日}

\begin{document}

\maketitle

\section*{第十二章}

\begin{enumerate}
    \item[12.1]
    天狼星的光度为$L = 25L_\odot$，则在太阳附近其能量密度为：
    \[
        u = \odv{E}{V} = \frac{L\d{t}}{4\pi D^2 \cdot c\d{t}} = \frac{L}{4\pi D^2 c} = 3.95 \times 10^{-15}\,\mathrm{erg \cdot cm^{-3}}
    \]
    另一方面，从书中的表(12.1)可得，恒星的背景光贡献的能量密度为：
    \[
        u_{\mathrm{star}} = 1.05 \times 10^{-12}\,\mathrm{erg \cdot cm^{-3}}
    \]
    则天狼星贡献的能量密度占总星光比例为：
    \[
        \frac{u}{u_{\mathrm{star}}} = \frac{3.95 \times 10^{-15}\,\mathrm{erg \cdot cm^{-3}}}{1.05 \times 10^{-12}\,\mathrm{erg \cdot cm^{-3}}} = 0.38\%
    \]

    \bigskip
    \item[12.2]
    首先计算能量范围在$10.0 < h\nu < 13.6\,\mathrm{eV}$的光子对应的波长范围：
    \[
        \lambda = \frac{c}{\nu} = \frac{hc}{h\nu}
    \]
    则有：
    \[
        10.0 < h\nu < 13.6\,\mathrm{eV} \implies 912\,\ang < \lambda < 1240\,\ang
    \]
    对于(12.7)式的能谱，由于等号右边是波长的函数，可以将等号左边的$\nu u_\nu$替换为$\lambda u_\lambda$，证明如下：
    \[
        \nu = \frac{c}{\lambda} \implies \d{\nu} = -\frac{c}{\lambda^2} \d{\lambda} = -\frac{\nu}{\lambda} \d{\lambda} \implies \lambda \d{\nu} = -\nu \d{\lambda}
    \]
    而单位光谱间隔的能量密度可以分别以频率和波长的形式表示为：
    \[
        u_\nu \d{\nu} = u_\lambda \d{\lambda}
    \]
    两式相除即得：
    \[
        \nu u_\nu = \lambda u_\lambda
    \]
    能量范围在$10.0 < h\nu < 13.6\,\mathrm{eV}$的光子跨越了(12.7)式中的两段能谱，则其数密度为：
    \[
        n = \int_{0.0912\,\um}^{0.124\,\um} \frac{u_\lambda}{hc/\lambda} \d{\lambda} = \frac{1}{hc} \int_{0.0912\,\um}^{0.124\,\um} \lambda u_\lambda \d{\lambda}
    \]
    代入(12.7)式得：（系数$k_1 = 1.287 \times 10^{-9}\,\mathrm{erg \cdot cm^{-3}}$，$k_2 = 6.825 \times 10^{-13}\,\mathrm{erg \cdot cm^{-3}}$）
    \begin{align*}
        n
        &= \frac{1}{hc} \ab[\int_{0.0912}^{0.11} k_1 (\lambda/\um)^{4.4172} \d{\lambda} + \int_{0.11}^{0.124} k_2 (\lambda/\um) \d{\lambda}] \\
        &= \frac{1}{hc} \ab[\int_{0.0912}^{0.11} (k_1 \times 1\,\um) \lambda_0^{4.4172} \d{\lambda_0} + \int_{0.11}^{0.124} (k_2 \times 1\,\um) \lambda_0 \d{\lambda_0}],\quad \lambda_0 = \lambda/\um \\
        &= \frac{1}{hc} \ab(\int_{0.0912}^{0.11} k'_1 \lambda_0^{4.4172} \d{\lambda_0} + \int_{0.11}^{0.124} k'_2 \lambda_0 \d{\lambda_0})
    \end{align*}
    其中$\lambda_0 = \lambda/\um$是无量纲的量，两个系数变为：
    \[
        k'_1 = 1.287 \times 10^{-13}\,\mathrm{erg \cdot cm^{-2}},\quad k'_2 = 6.825 \times 10^{-17}\,\mathrm{erg \cdot cm^{-2}}
    \]
    计算可得：
    \begin{align*}
        n
        &= \frac{1}{hc} \times (9.715 \times 10^{-20}\,\mathrm{erg \cdot cm^{-2}} + 1.118 \times 10^{-19}\,\mathrm{erg \cdot cm^{-2}}) \\
        &= \frac{2.089 \times 10^{-19}\,\mathrm{erg \cdot cm^{-2}}}{1.986 \times 10^{-16}\,\mathrm{erg \cdot cm}} \\
        &= 1.05 \times 10^{-3}\,\mathrm{cm^{-3}}
    \end{align*}

\end{enumerate}

\end{document}
